\section{Controlled removal of the middle homology}\label{ct:section}

Suppose, towards a contradiction, that the group constructed above is the
fundamental group of a closed aspherical topological four-manifold.  We use
the marked systems supplied by Proposition~\ref{st:periodic-data}.  Thus,
after finitely many stabilizations on the compact quotient, the ordinary
reflection development is a manifold
\[
                 Y=\bigcup_{v\in C}D_v.
\]
Here $D_v$ is a stabilized copy of the chamber, its boundary is $S$, and
its group is the free group $H$.  Boundary identifications are understood
with the signs induced by the orientation of $Y$; signs in hyperbolic
planes are absorbed by changing one basis vector.  There are two
labelled systems in $Y$:
the chamber-wise immersed hyperbolic bases, denoted $\mathcal A$, and
disjoint standard hyperbolic pairs, denoted $\mathcal B$.  Each pair of
$\mathcal B$ has an embedded neighborhood homeomorphic to
$(S^2\times S^2)\setminus\operatorname{int}D^4$.  Corresponding spheres in
the two systems are based-homotopic.  All these data, including the
homotopies, are lifts of finitely many compact data on the quotient.

We will remove an expanding finite part of $\mathcal B$ and use the
remaining part of $\mathcal A$ to perform surgery.  The latter operation
requires geometric duals.  We obtain them by applying the proof of the
disc embedding theorem to a specially constructed null-loop input;
we do not assume that $H$ is a good group.

There are three tasks. First we cap the noncentral boundary factors and
remove the inner standard pairs, leaving a compact manifold with the
original boundary marking. Next we prepare the remaining sphere
systems and place their upper grope branches over contractible
direction patches. This supplies the null-loop input. Finally the
tower construction and sphere surgeries remove the remaining middle
homology and produce a graph-type filling.

\subsection{Complementary fillings and the peeling operation}

Write $r=\operatorname{rank}H$.  A stabilization of $D$ has homotopy type
\[
        \bigvee^r S^1\vee\bigvee^{2k}S^2
\]
for some $k$, and its middle module has the specified hyperbolic basis.
Neither the boundary nor its map to $BH$ changes under stabilization.

\begin{lemma}[Complementary fillings]\label{ct:fillings}
For every such stabilized chamber there is a simply connected compact
oriented topological four-manifold $J$ with boundary $-S$ such that
\[
 H_2(S;\Z)\longrightarrow H_2(J;\Z)
 \quad\text{is an isomorphism},\qquad H_3(J;\Z)=0.
\]
Moreover, $D\cup_S J$ is simply connected, and inclusion of $D$ identifies
its second homology and intersection form with those of $D\cup_S J$.
\end{lemma}

\begin{proof}
First recall the consequences of the graph-type Poincar\'e pair $(P,S)$
recorded in Corollary~\ref{ch:boundary-homology}.  Duality and its
ordinary homology sequence give
\begin{equation}\label{ct:boundary-homology}
 H_1(S;\Z)\cong H_1(P;\Z)\cong\Z^r,
 \qquad H_2(S;\Z)\cong H^1(P;\Z)\cong\Z^r.
\end{equation}
The same corollary gives the surjectivity of $\pi_1(S)\to H$.

Start with the oppositely oriented copy of $D$.  Represent a free basis
of its fundamental group by disjoint embedded circles.  Their normal
bundles are trivial, and general position makes the circles disjoint
from the finite sphere system.  Interior surgery on these circles kills
the free generators.  The old sphere classes retain their hyperbolic
intersection and framing data after augmentation to $\Z$.  In the now
simply connected manifold, the sphere embedding theorem with geometric
duals realizes a Lagrangian by disjoint framed embedded spheres.  Surgery
on these spheres preserves simple connectivity and removes the old
hyperbolic planes.  This use of the disc theorem is in the simply
connected case only; see \cite[Theorem~B]{PRT}.

Let $J$ be the result.  Each circle surgery increases Euler characteristic
by two, and each sphere surgery decreases it by two.  Hence
\[
 \chi(J)=(1-r+2k)+2r-2k=1+r.
\]
The boundary is connected and $J$ is simply connected.  Thus
$H^1(J,S;\Z)=0$ by the cohomology sequence of the pair, and
Poincar\'e--Lefschetz duality gives $H_3(J;\Z)=0$.  Since $J$ has nonempty
boundary, $H_4(J;\Z)=0$, and therefore $b_2(J)=r$.

There is no integral index ambiguity in this calculation.  The relevant
part of the duality sequence is
\begin{equation}\label{ct:cap-sequence}
 0\longrightarrow H_2(S;\Z)\longrightarrow H_2(J;\Z)
 \longrightarrow H^2(J;\Z)\longrightarrow H_1(S;\Z)
 \longrightarrow0.
\end{equation}
Here $H_3(J,S;\Z)=H^1(J;\Z)=0$ accounts for the initial zero.
The universal coefficient theorem makes $H^2(J;\Z)$ free, because
$H_1(J;\Z)=0$.  Every torsion element of $H_2(J;\Z)$ would consequently
belong to the image of the injected free group $H_2(S;\Z)$, so there is
no such element.  The first two nonzero groups in
\eqref{ct:cap-sequence} now have the same rank $r$.  The middle map has
rank zero and takes values in a free group, hence is zero.  Exactness
proves that the first map is an integral isomorphism.

The map $H_2(S;\Z)\to H_2(D;\Z)$ is zero: the next map
$H_2(D;\Z)\to H_2(D,S;\Z)$ is the nonsingular hyperbolic intersection
map.  Also $H_1(S;\Z)\to H_1(D;\Z)$ is the isomorphism in
\eqref{ct:boundary-homology}.  Mayer--Vietoris, together with the
isomorphism just proved for $J$, therefore identifies
$H_2(D\cup_S J;\Z)$ with $H_2(D;\Z)$.  All products are unchanged,
since their representatives lie in the interior of $D$.
Finally, the surjectivity of $\pi_1(S)\to H$ and van Kampen show that
$D\cup_S J$ is simply connected.
\end{proof}

Let $|v|$ denote Coxeter word length.  Lemma~\ref{rf:shelling} identifies
the finite chamber union
\[
                  U_L=\bigcup_{|v|\leq L}D_v
\]
with an iterated boundary connected sum, with the natural markings on
its boundary factors.  Fill every noncentral boundary factor by the
corresponding $J_v$ of Lemma~\ref{ct:fillings}.  The resulting manifold
$U_L^+$ has boundary $S$ and admits the description
\begin{equation}\label{ct:filled-sum}
 U_L^+\cong D_e\mathbin{\#}
       \mathop{\#}_{0<|v|\leq L}(D_v\cup_S J_v).
\end{equation}
To see this description, retain the connecting three-balls in the
boundary-sum construction.  Gluing a punctured filling to a noncentral
factor closes that factor away from its connecting ball and turns the
connecting boundary neck into an interior connected-sum neck.  Repeating
this one factor at a time gives \eqref{ct:filled-sum}.

It follows that $\pi_1(U_L^+)=H$, with the central boundary marking.
On $\pi_1(U_L)$ the inclusion is the projection of the free product of
the chamber groups onto its central factor.  The universal cover of
$U_L^+$ has reduced homology only in degree two, freely based over
$\Z H$ by all the chamber sphere systems, including stabilizations.
Its intersection form is their orthogonal hyperbolic sum.  These facts
also follow directly by lifting the connected-sum necks in
\eqref{ct:filled-sum}; all the closed noncentral factors are simply
connected and have precisely the middle classes computed above.

For an integer $R$, replace every standard neighborhood in $\mathcal B$
whose label belongs to a chamber $|v|\leq R$ by a four-ball.  Call this
operation \emph{peeling}.  The common boundary is a standard $S^3$.
There is a continuous pinch map from each old neighborhood to the new
ball, fixed on that boundary, and these maps combine with the identity
outside to a global pinch map.  Both pieces and the common boundary are
simply connected, so peeling preserves the fundamental group.  The
relative homology sequence deletes exactly the hyperbolic pair carried
by that neighborhood and introduces no third homology.

We first perform this operation in the unbounded $Y$, obtaining $Y_R$.
For $L$ large enough to contain the selected neighborhoods and all their
marking homotopies, it also defines a compact manifold
\[
                        Z=Z_{R,L}
\]
by peeling $U_L^+$.  The order of filling and peeling is immaterial
because the standard neighborhoods are in the interior.  The included
homotopies identify the removed summands with exactly the chamber pairs
labelled by $|v|\leq R$.  We have therefore proved
\begin{equation}\label{ct:Z-homology}
 \partial Z=S,\qquad \pi_1 Z=H,\qquad
 \widetilde H_i(\widetilde Z;\Z)=0\quad(i\ne2),
\end{equation}
and $H_2(\widetilde Z;\Z)$ is freely based by the projected chamber
pairs with $R<|v|\leq L$, with their orthogonal hyperbolic form.

\subsection{Uniform finite preparation}

A compact object is \emph{carried within $b$ galleries of $v$} if its
image is contained in chambers whose labels are within word distance
$b$ of $v$.  Replacement balls are assigned the carriers of their old
standard neighborhoods.  We use this terminology for sphere maps,
paths, homotopies and grope systems, including all their caps.  A
uniform propagation bound is one value of $b$ valid for the whole
labelled family.  It is stronger than a bound after an unspecified
permutation of the labels.

\begin{lemma}[Finite preparation]\label{ct:preparation}
Fix a positive integer $h$.  In $Y_R$, the projected pairs of
$\mathcal A$ labelled by $|v|>R$ admit the following preparation, with
propagation bounded independently of $R$:
\begin{enumerate}
\item The two halves are represented by framed immersed spheres
      $a_i,b_i$ with $a_i\cap b_j$ consisting of the designated single
      point for $i=j$ and empty for $i\ne j$.
\item The remaining intersections among the $a_i$ are paired by
      framed Whitney discs.  Each Whitney disc is replaced by a
      properly immersed capped grope of height at least $h$ with the
      same framed attaching region.
\item These gropes have the Clifford dual capped surfaces used in
      \cite[Definition~3.1 and Lemma~4.1]{PRT}.  In particular, their
      caps miss the \emph{entire} dual capped surfaces, and the bodies
      of the latter lie in a boundary collar of the complement of the
      $a_i$.
\end{enumerate}
The construction uses no complementary filling $J_v$.  Beyond a fixed
gallery distance from the peeling region, it can be carried out
separately inside individual chambers.
\end{lemma}

\begin{proof}
First project the outer immersed systems through the pinch and put them
in general position.  The deleted and retained blocks were orthogonal,
since their classes agree with the correspondingly labelled standard
pairs.  The retained forms and quadratic self-intersection data are
therefore still hyperbolic over $\Z\pi_1(Y_R)$.  Choose framed immersed
representatives and cancel the algebraic intersection pairs.  The
framing corrections and, where necessary, local cusp changes can be
chosen from finitely many local drawings.  The group in question is a
free product of free chamber groups, hence is torsion-free; there is no
order-two ambiguity in these quadratic data.

We justify carefully the uniformity of the Whitney discs required at
this and subsequent finite stages.  The initial lifted data have
finitely many types under a cocompact group of translations of $Y$.
A gallery ball of fixed radius contains a uniformly bounded number of
chambers and of pieces of those data.  There are thus finitely many
local templates, even after specifying which of the standard
neighborhoods meeting that ball have been peeled.  Paths on participating
spheres are taken inside their bounded carriers.  Common long whiskers
conjugate an intersection label but need not be included in the
geometric Whitney circuit.

If a circuit enters a replacement ball, move its portions there to
chosen arcs on the boundary three-sphere.  This gives a circuit in the
unpeeled space.  Nullity is unchanged, since both the deleted
punctured $S^2\times S^2$ and the replacement ball are simply connected
and have simply connected boundary.  For every null circuit among the
finitely many resulting templates, choose one compact filling in $Y$.
There is a finite maximum of their propagation bounds.  Translate the
chosen filling to the location of the circuit and apply the actual
global pinch map.  This constructs a filling in $Y_R$.

A filling may encounter additional peeled neighborhoods.  No further
choice of fillings is then needed: the global pinch is already defined
on them, and each such neighborhood has the same bounded carrier.
Consequently this procedure is uniform in $R$.  It is a finite-template
argument for the specified circuits, not a bound for arbitrary null
loops.  Repeating a fixed finite number of stages preserves the
argument: enlarge the carrier radius and refine the finite template
list at each stage.

Now apply the geometric Casson lemma to the two halves of the sphere
system, removing the paired mixed intersections by regular homotopies.
There is no iterative repair cascade here.  Use the local construction
of \cite[Lemma~3.3]{PRT} with the two surface arguments taken to be the
entire possibly disconnected families.  Preselect all framed pairing
Whitney discs with disjoint boundaries; their interiors miss the other
Whitney boundary arcs by $2+1<4$ general position.  First push the first
family off all these interiors by fingers along arcs in the Whitney
discs.  Choose the arcs away from their finitely many intersections
with other Whitney discs and the second family.  The finger supports
then leave the other Whitney discs and attaching boundaries fixed.
Next Whitney-move the second family using the same discs.  The new
intersections in each phase are only within that phase's family;
intersections between Whitney interiors cause only allowed
second-family intersections.  Thus the whole batch stays in the
original bounded Whitney neighborhoods.  No later disc is transported
or chosen again through an expanding sequence of carriers.  Local
finiteness permits these bounded batches on the unbounded $Y_R$.
This produces the geometric duality in (1), while permitting
self-intersections within either half.  The remaining double points
of the first half admit framed Whitney discs.  Tube their interiors
off the first-half spheres using the geometric duals.  All the
Whitney circuits involved are the null circuits just discussed.

For completeness, we verify the algebraic input to the grope preparation
in the complement $M_0$ of small plumbed neighborhoods of the $a_i$.
The punctured geometric dual $b_i$ bounds an $a_i$-meridian in $M_0$.
Since these meridians normally generate the kernel, inclusion induces
$\pi_1M_0\cong\pi_1Y_R$.  Thus the following calculation retains all
intersection labels in the required complement group ring.

At one point of every paired double point take its Clifford torus.
Cap its two basis curves by standard meridional discs and tube each
meridional disc off its incident $a_i$-sheet into a punctured parallel
of $b_i$, with the product framing.  Denote these caps by $C_\ell$,
and let $i(\ell)$ be the index of the dual used.  Their equivariant
intersection calculation is the Clifford-cap calculation in the proof
of \cite[Section~5]{PRT}.  The local meridional pieces have disjoint
interiors.  Choose the thin tubing tracks away from the finitely many
other designated dual points; local changes of tracks introduce only
oppositely signed points with the same group label.  The remaining
intersections are those between the corresponding based parallels of
$b_{i(\ell)}$.  Hence their sums are obtained from the pairings of the
$b_i$ by the basing units and orientation signs.  Likewise, a cap's
self-intersection sum is that of its framed $b_i$ parallel, together
with cancelling local terms.  Since the mixed Casson regular homotopies
preserved $\lambda(b_i,b_j)=0$ and $\mu(b_i)=0$, this proves
\[
 \lambda(C_\ell,C_m)=0,\qquad \mu(C_\ell)=0
 \quad\text{for every }\ell,m,
\]
including distinct parallel caps using the same dual sphere.

The Clifford bodies lie in a boundary collar of $M_0$.  Following
\cite[Section~5]{PRT}, contract temporary parallel copies and push the
Whitney discs off those contractions before using them as geometric
dual spheres.  At the remaining intersections with Whitney discs,
tube the Clifford caps into parallel copies of these temporary duals.
Lemma~4.1 of \cite{PRT}
now applies with its full cap-pairing hypothesis.  Its two separation
batches and height raising provide the prepared Clifford copies,
whose entire capped surfaces are geometrically dual in the sense of
\cite[Definition~3.1]{PRT}.  These prepared copies, including all their
caps, are the protected surfaces used from this point onward.
The preparation in \cite[Lemma~4.1]{PRT} produces height-two Whitney
gropes, separates their caps from the protected Clifford caps, and
raises the height to $h$. The collar levels and the full dual-cap
disjointness are part of this preparation, not an assertion made after
the caps have been changed.
There are exactly two cap-family separation batches before height
raising.  Each uses the same no-new-mixed-pairs invariant, with its
protected bodies at fixed collar levels.

Each of the operations before height raising involves only a fixed
number of kinds of tubing, push-off, Whitney-disc choice and regular
homotopy.  Their simultaneous use does not require traversing the
whole list of sphere components.  At each fixed stage, bounded
carriers and finite drawings bound the number of other carriers
meeting a given one.  Divide the operations into finitely many
classes with disjoint enlarged carriers and perform one class at a
time.  This gives a number of local stages independent of $R$.
New null circuits are handled by the finite-template procedure above.

For height raising there is also an explicit size estimate:
\cite[Section~2.8.3]{FQ}, in the January 2013 numbering, bounds the
increase of component-image diameter by a factor of seven at each
step.  A fixed requested height therefore has a fixed propagation
bound.  One can apply this estimate to the gallery map into the
Davis complex: extend that map over each replacement ball by coning
its boundary image inside a metric ball of uniformly bounded radius.
This is possible by CAT(0) convexity.  A point in a chamber is mapped
within a fixed distance of its label, and properness of the Coxeter
action turns a fixed metric bound into a fixed word-distance bound.
The construction uses parallel transverse gropes and
neighborhood operations; it preserves the protected Clifford capped
surfaces.  All these operations can be made locally finite in the
unbounded development.

Finally, outside a fixed buffer of the peeling region the input is
the original system in an individual chamber.  Its intersection
cancellations hold already over that chamber's group ring.  The same
preparation, including the null discs, can consequently be chosen in
that chamber.  Only a bounded buffer of the switch between the
peeled and unpeeled data requires the preceding translated templates.
This proves the last assertion and permits later truncation without
cutting off an unfinished preparation.
\end{proof}

\subsection{Directions at infinity}

Let $\Sigma_C$ be the ordinary Davis complex and let $o$ be its vertex
labelled by $e\in C$.  The gallery map of Lemma~\ref{rf:davis-map} sends
the inner part of $D_v$ to $v$, interpolating across the dual collars.
The link of $o$ is $S$.  A small metric sphere of radius $t$ about $o$
is consequently identified with $S$ by radial scaling in its conical
star.  Write $\rho$ for radial projection onto this sphere, defined
outside the open ball of radius $t$.

\begin{lemma}[Small directions and the central marking]\label{ct:directions}
Fix a propagation bound $b$.  For objects carried within $b$ galleries
of labels $v$ with $|v|\to\infty$, their direction images in $S$ have
diameter tending uniformly to zero.  This assertion continues to hold
after peeling for the whole region of bounded propagation around the
outer systems, if $R$ is sufficiently large.

Let $\beta:S\to BH$ be the central boundary marking.  On that region,
the homomorphism induced by the direction map followed by $\beta$
agrees, up to the common change of base path, with inclusion into
$Z$ followed by its classifying map.  In particular, every loop in an
image system whose entire direction image lies in a contractible
patch of $S$ is nullhomotopic in $Z$.
\end{lemma}

\begin{proof}
Here is the radial estimate that supplies the first assertion.  If
$x,y\in\Sigma_C$ have distance at least $s>t$ from $o$, put them first
on the same metric sphere, of radius the smaller of their two radii.
The distance between the resulting points is at most $2d(x,y)$,
by the triangle inequality and
$|d(o,x)-d(o,y)|\leq d(x,y)$.  CAT(0) comparison on the two radial
segments then gives
\begin{equation}\label{ct:radial-estimate}
                  d(\rho(x),\rho(y))
                  \leq \frac{2t}{s}\,d(x,y).
\end{equation}
The gallery map takes each chamber into a set of uniformly bounded
diameter.  Thus a bounded gallery carrier has bounded image diameter.
Properness of the Davis complex, together with its cocompact Coxeter
action, makes its distance from $o$ tend to infinity with the word
length of its center.  Equation~\eqref{ct:radial-estimate} applies.

We only need the direction map outside the central core.  Consider a
peeled standard neighborhood that meets an outer object or any of
its fixed enlarged carriers.  Both the object and the whole standard
neighborhood have bounded propagation.  The label of this neighborhood
therefore has length at least $R$ minus a fixed constant.  Its entire
boundary three-sphere, not merely the arcs encountered by the object,
has direction image of arbitrarily small diameter for large $R$.
Choose a slightly larger coordinate ball of $S$ containing this image.
Contracting that ball extends the boundary map across the replacement
four-ball, with image still in the chosen ball.  These extensions
agree with the unchanged direction map on all their boundaries.

The near-central peeled neighborhoods need not admit such small
extensions.  They are disjoint from all the carriers under consideration
once $R$ is large, and no direction map on them is used.  This proves
the asserted far-out extension and its uniform smallness.

Before peeling, the composition with $\beta$ extends across the
central chamber by its given map to $BH$.  Its restriction to the
central boundary agrees with that map: this is the boundary-direction
identification in Lemma~\ref{rf:davis-map}.  In a noncentral chamber
the gallery map collapses its inner part to the noncentral vertex,
so every loop of its chamber group has constant direction there.
The shelling description makes the chamber groups the free factors
of the fundamental group of a finite ball of chambers.  The induced
homomorphism is thus exactly the projection retaining the central
factor.  By \eqref{ct:filled-sum} this is also the homomorphism induced
by passage to $U_L^+$.

Replacing a simply connected piece by a simply connected ball does
not alter this calculation.  Loops through the replacement can be
moved to the boundary of the piece, and the direction extension
agrees there with the old map.  The comparison therefore holds on
the required region in $Z$.  If a whole image system maps into a
contractible patch, every one of its loops, including loops passing
through arbitrarily many intersections, has trivial image under
$\beta$.  Since $BH$ is a classifying space and $\pi_1Z=H$, these
loops are nullhomotopic in $Z$.
\end{proof}

\subsection{Upper grope separation}

We state explicitly how the controlled data supply the input used
after the good-group step of the disc embedding proof.  The units to
be separated are the upper capped subgropes attached to individual
tips of the retained stages.  Separating only whole Whitney gropes
would not suffice to make the future cap discs mutually disjoint.

\begin{lemma}[Null upper replacements]\label{ct:upper-replacement}
Suppose a finite prepared Whitney-grope system as in
Lemma~\ref{ct:preparation} has at least eight surface stages.  Retain
its first five stages.  Suppose the upper capped subgropes, of height
at least three, have uniformly small direction images in the region
of Lemma~\ref{ct:directions}.  If their diameters are sufficiently
small, depending only on a fixed finite coordinate cover of $S$, the
upper subgropes can be replaced by mutually disjoint framed immersed
cap discs.  Their self-intersection loops are nullhomotopic in $Z$.
The retained stages and the entire protected Clifford capped
surfaces remain unchanged.
\end{lemma}

\begin{proof}
We first specify the relative setting. Remove a sufficiently small
regular neighborhood of the five retained stages, truncating each upper
attaching collar at its boundary. The upper components are now properly
immersed disk-like capped gropes of height at least three. By the duality
condition in Lemma~\ref{ct:preparation}, their entire images miss the
protected Clifford capped surfaces. Choose a neighborhood of the upper
system disjoint from those surfaces and from the retained stages except
at the prescribed attaching regions. All the following modifications
are neighborhood operations in this relative manifold. They therefore
leave the lower stages, attaching framings, and whole Clifford systems
unchanged, including the unused boundary-collar levels needed later.

We use the following precise content of Freedman--Quinn's controlled
separation lemma \cite[Lemma~2.10.1]{FQ}. For a map from a four-manifold
to a locally compact metric space, a properly immersed union of
disk-like capped gropes of height at least three, with sufficiently
small component images, admits new caps for the \emph{same bodies}
such that distinct components are disjoint and their images remain
small. Its compact-case proof separates finitely many cover collections
before separating the members of each collection
\cite[Section~2.10.2]{FQ}. We retain this two-phase construction because
we need a bound on gallery propagation only in its first phase.

Choose compact sets $K_1,\ldots,K_q$ covering $S$, each contained in an
open coordinate ball $V_j$. Choose $\epsilon>0$ so that the closed
$\epsilon$-neighborhood of every $K_j$ lies in $V_j$. Assign each upper
component to a $K_j$ met by its direction image. If every initial image
has diameter less than $\delta$, that component lies within distance
$\delta$ of its assigned $K_j$.

For the first phase, regard all components assigned to $K_j$ as one
collection. The construction in \cite[Section~2.10.2]{FQ} raises height
by one, contracts the added stages of the first collection and pushes
the later collections off, then continues through the collections in
order. The original bodies are retained. Its component-diameter estimate is
$3^{q-1}7^2\delta$. A point of each retained body stays within
$\delta$ of $K_j$, so the entire resulting $j$th collection lies in
the $3^{q+3}\delta$-neighborhood of $K_j$. Choose $\delta$ with
$3^{q+3}\delta<\epsilon$, leaving room for a small regular-neighborhood
thickening. The number of height-raising and contraction passes depends
only on $q$. The same size estimates apply to the gallery map into
$\Sigma_C$, extended over the replacement balls as in
Lemma~\ref{ct:preparation}. Since each original body is retained, its
image anchors the resulting component within a bounded metric distance
of the original label. Properness of the Coxeter action converts this
into a gallery bound independent of the number of components and of $R$.

The collections now have pairwise disjoint images. Choose disjoint
regular neighborhoods of their entire images, with the $j$th
neighborhood still mapping into $V_j$. Inside each such neighborhood,
apply the collection-separation procedure again, this time taking
each individual upper component as a collection. There may be
arbitrarily many of them. The height is raised once more inside that
neighborhood before contracting the added stages, so this second
phase again retains the original bodies. No estimate depending on
their number is needed: all operations stay in the already chosen
neighborhood. The resulting upper components have pairwise disjoint
neighborhoods. Totally contract each in its own neighborhood to obtain
the claimed mutually disjoint framed immersed cap discs.

Every loop in a resulting cap image lies in its color neighborhood,
whose entire direction image is contained in the contractible $V_j$.
Lemma~\ref{ct:directions} therefore makes that loop nullhomotopic in
$Z$. The second phase need not bound the gallery diameter of an
individual cap. It preserves instead the containment of its whole
image in a direction patch, which is exactly the condition needed
for nullity.
\end{proof}

\subsection{Choice of parameters and the tower construction}

Fix the coordinate cover in Lemma~\ref{ct:upper-replacement}. Request
height eight in Lemma~\ref{ct:preparation}. Its uniform propagation
bound, together with the bounds for the standard neighborhoods and
marking homotopies, controls the entire preparation near the peeling
boundary. The first phase of upper separation adds only the fixed
height extension and $q$ collection passes described above. Let $b$
bound the gallery propagation of all these operations, independently
of $R$. Choose $R$ so large that the direction estimates of
Lemma~\ref{ct:directions}, including the replacement-ball extensions,
meet the chosen cover margins for every such carrier.

We now truncate the prepared family, before carrying out the
collection-separation operations. Choose $L$ large enough to contain
every peeled standard neighborhood and its marking homotopies, and
the whole $b$-buffer around the transition between peeled and
unpeeled data. Beyond that buffer, Lemma~\ref{ct:preparation} constructs
the sphere and Whitney-grope data entirely inside their own chambers,
using their chamber group-ring pairings. Retain precisely the planes
whose labels satisfy $R<|v|\le L$, together with their complete
preparations. The preparations at the outer boundary lie in chamber
interiors, so none has a partner or an attaching region outside
$U_L$. This gives a finite family in the peeled copy of $U_L$, before
any complementary filling is used.

Cap its noncentral boundary factors to form $Z=Z_{R,L}$. Apply the
first phase of
Lemma~\ref{ct:upper-replacement} to this finite family. All of its
operations can be made in arbitrarily small neighborhoods of the
retained upper system inside that peeled region; the bound $b$ supplies
the direction
margins independently of the truncation. Apply the second phase inside
the resulting disjoint color neighborhoods. Neither phase requires
an enlargement of $L$ or uses the complementary fillings. The second
phase supplies the immersed null-loop caps without a uniform gallery
bound on their individual images.

Write
\[
                  M'=Z\setminus\bigcup_i\operatorname{int}\nu(a_i)
\]
for the complement used in the Whitney construction.  The spheres
$a_i$ may still be immersed; the notation denotes their standard
plumbed neighborhoods.  The prepared geometric dual spheres imply
\begin{equation}\label{ct:complement-group}
                   \pi_1(M')\xrightarrow{\ \cong\ }\pi_1(Z).
\end{equation}
Indeed the kernel of the inclusion map is normally generated by
meridians of the $a_i$, and the punctured dual spheres bound these
meridians in the complement.  General position supplies surjectivity.
This is the usual $\pi_1$-negligibility provided by geometric duality.

The new caps therefore have null double-point loops in $M'$, not
just in $Z$.  They are mutually disjoint, although each may have
self-intersections. Direction control has now supplied the needed
fundamental-group condition. The remaining constructions require that
condition and the protected geometric duals, but no bound on the
diameter of the null homotopies. Symmetrically contract the fifth retained
surface stage.  The result is a height-four capped Whitney grope
system, still with mutually disjoint caps.  Its cap
self-intersections occur in algebraically cancelling pairs, and
their loops are parallel copies of the previous null loops.  This
is exactly the intermediate input in the proof of
\cite[Lemma~4.2]{PRT}, after that proof has used goodness.  The
retained Clifford capped surfaces are still geometrically dual in
the stronger sense of \cite[Definition~3.1]{PRT}.

We explain why the remainder of that proof applies without any
assumption on $H$.  Cap the null double-point loops by immersed
discs in $M'$.  The bodies of the Clifford capped surfaces lie at
the protected boundary-collar levels.  The attaching loops of the
new discs miss that collar, so the null homotopies can be chosen
to miss the Clifford \emph{bodies} by pushing them out of the
collar.  They may meet the Clifford caps and may meet the Whitney
gropes arbitrarily.  Correct their boundary framings, push their
intersections with the gropes down, and tube into geometrically
dual spheres obtained by contracting suitable parallel copies of
the Clifford systems.  These operations produce the tower caps.
At the next collar level, contract the remaining Clifford capped
surfaces and push the tower caps off those contractions.  The
resulting spheres are geometric duals to the one-storey capped
towers.  These are the successive operations in the remainder of
the proof of \cite[Lemma~4.2]{PRT}; no subsequent step uses goodness.

One-storey capped towers with four surface stages contain the
embedded discs required for the Whitney moves, with the same
framed attaching regions.  Applying the final steps of
\cite[Section~5]{PRT}, first move the original sphere duals off the
towers by tubing into the tower duals.  Then use the embedded
Whitney discs to remove the intersections of the $a_i$.  We obtain
disjoint framed embedded representatives of their original
Lagrangian classes, together with geometric transverse spheres.
The changes of upper caps do not change the retained Whitney
attaching data or the corrected sphere classes.  In particular,
the symmetric-contraction independence of cap choices in
\cite[Lemma~3.4]{PRT} applies at the retained stages.

There is deliberately no assertion that the null homotopies, the
tower caps, or the final embedded Whitney discs have small gallery
diameter.  From \eqref{ct:complement-group} onwards they may range
throughout $Z$.  Only their existence, their relative attaching
data, and the geometric duals are needed.

\Needspace{7\baselineskip}
\begin{proposition}[Removal of the remaining planes]\label{ct:removal}
The compact manifold $Z$ admits interior surgeries producing a
compact oriented topological four-manifold $V$ with boundary $S$,
fundamental group $H$, and contractible universal cover.  The
boundary map to $BH$ is the original central marking.
\end{proposition}

\begin{proof}
Perform surgery on the disjoint framed Lagrangian spheres just
constructed.  Their geometric duals kill the meridians in their
complement, so the surgeries preserve the fundamental group.
All changes are in the interior; the boundary and its marking are
unchanged.

For clarity, let there be $m$ remaining planes and let $W$ be the
five-dimensional surgery trace.  With regular $\Z H$ coefficients,
the relative chain complex of $(W,Z)$ is concentrated in degree
three, with module $(\Z H)^m$.  Its boundary map to
$H_2(Z;\Z H)$ sends the basis to the Lagrangian basis.  By
\eqref{ct:Z-homology} this is a split injection.  Consequently
\[
 H_3(W;\Z H)=0,\qquad
 H_2(W;\Z H)\cong(\Z H)^m,
\]
the latter module represented by the dual half, and the other
positive-degree homology groups vanish.  Read the same trace from
$V$.  Its relative chain complex is concentrated in degree two,
again with module $(\Z H)^m$.  The map
\[
            H_2(W;\Z H)\longrightarrow H_2(W,V;\Z H)
\]
is the pairing of the surviving dual classes with the surgered
Lagrangian, hence is an isomorphism.  The sequence of $(W,V)$
therefore gives $H_i(V;\Z H)=0$ for every $i>0$.

The universal cover of $V$ is simply connected and acyclic. A
topological manifold has CW homotopy type~\cite[Corollary~1]{MilnorCW};
in the boundary case one may first append an exterior open collar,
which does not change homotopy type. The homological Whitehead theorem
therefore makes this universal cover contractible
\cite[Corollary~4.33]{Hatcher}. Thus $V$ is a $K(H,1)$ and is homotopy equivalent
to the finite graph $BH$.  Under this equivalence its boundary
map is the retained central marking, as asserted.
\end{proof}

The manifold $V$ is the marked graph-type filling excluded by
Proposition~\ref{ch:chamber}.  This contradiction rules out the hypothetical
closed aspherical topological four-manifold and completes the proof
of Theorem~\ref{thm:main}.
